% §1.5 Why "Atom"?
%
% Cross-reference resolution status:
%
% ALREADY DEFINED (within chapter 1, by this section):
%   \label{sec:why-atom}                    -- this section
%
% RESOLVED:
%   \Cref{sec:three-faces}, \Cref{sec:quantifying-noise},
%   \Cref{sec:boolean-universe-first-look}, \Cref{sec:sampling-rates} -- ch.1
%   \Cref{ch:binary-channels}, \Cref{sec:model-order}                 -- ch.2
%   \Cref{ch:boolean-algebra}, \Cref{sec:composition-theorem},
%   \Cref{sec:miller-rabin}                                           -- ch.3
%   \Cref{ch:axioms-kronecker}, \Cref{sec:axiom-1}, \Cref{sec:axiom-2},
%   \Cref{sec:kronecker}, \Cref{sec:bernoulli-bool-formal}            -- ch.4
%
% Monograph register: justifies "atom" (two senses of atomicity) and
% states the trajectory of Part I and the rest of the book compactly. The
% formal type name first appears in the chapter-4 preview sentence.

\section{Why ``Atom''?}
\label{sec:why-atom}

The noisy bit is atomic in two senses. First, it has no internal Bernoulli
structure to decompose: the input is one Boolean value, the output is one
Boolean value, and the entire behavior is two numbers, $\fprate$ and $\fnrate$.
A smaller part would require a sub-Boolean input or output, and there is
nothing below a Boolean. The atom is where internal accounting stops and
external accounting begins, and every larger object in the book, chained
channels, noisy logic gates, randomized algorithms, approximate sets, general
algebraic types, is assembled from noisy bits by a specific composition rule.
Second, the noisy bit sits at the minimum of the model-order\index{model order} ladder that
\Cref{sec:model-order} defines, the number of distinct error regimes an input
domain splits into. A general noisy bit has model order two, the blocks
$\{\top\}$ and $\{\bot\}$ carrying independent rates; the symmetric case
collapses to order one. No noisy system over a Boolean domain has a smaller
non-trivial channel matrix.

The claim is a working hypothesis until discharged, and chapters 2 through 4 do
so in three directions. \Cref{ch:binary-channels} gives the noisy bit a
$2 \times 2$ row-stochastic matrix $Q$, places the symmetric, asymmetric, and
Z-channel cases on the model-order ladder, and shows the atom closed under
serial composition (the matrix product of row-stochastic matrices is
row-stochastic). \Cref{ch:boolean-algebra} asks what AND, OR, and NOT do to
noisy bits. Gate by gate the answer is closed-form, but it is not again a noisy
bit: a gate over two distinct latent bits has a single rate on one side and up
to three on the other, so a pair of rates does not propagate through a Boolean
tree. What does propagate exactly is the error at each latent assignment of the
inputs; a single pair survives only as a prior average or as a worst-case bound
(\Cref{sec:boolean-formulas}), and corner products such as the composition
theorem (\Cref{sec:composition-theorem}) are exact only in their stated
regimes. One round of the Miller-Rabin\index{Miller--Rabin test} primality test
is a noisy bit with $\fnrate = 0$ and $\fprate \le 1/4$, and because every
round reads the same latent value the $k$-round bound $\fprate^{(k)} \le
(1/4)^k$ is the all-negative product rule rather than a separate analysis
(\Cref{sec:miller-rabin}). \Cref{ch:axioms-kronecker}
names the two independence assumptions used silently until then, element-wise
independence (\Cref{sec:axiom-1}) and conditional independence of per-block
rates (\Cref{sec:axiom-2}); proves the Kronecker factorization\index{Kronecker factorization} they purchase
(\Cref{sec:kronecker}), which is what makes the per-element description
sufficient; and writes the formal definition $\Bernoulli{\BoolType} = (X, Y,
Q)$ (\Cref{sec:bernoulli-bool-formal}), of which every atom in chapters 1 through 3
is an instance\index{Bernoulli Boolean}.

Beyond Part~I the atom is held fixed and its surroundings vary. Part~II keeps
the Boolean codomain and arranges many noisy bits into a structure: a set $S$
over a universe induces a membership query per element, an approximate set
makes each query a noisy bit (the Bloom filter being the canonical one-sided
instance, $\fnrate = 0$ with tunable $\fprate$), and the development is the
algebra of how set operations propagate error, how classification measures
quantify quality, and how much space a target accuracy demands. Part~III
replaces the Boolean codomain with an arbitrary algebraic type, lifting the
atom to approximate maps, relations, and structured values, where the
two-element model order ceases to be a ceiling. The atom is the unit that gets
composed; everything that follows is composition, and an honest account of what
composition preserves and what it does not.
